% !TeX root = ../main.tex
\section{Lambda calculus}
\label{sec:lambda}

This section and the next go beyond the slides; they give the formal basis
for the rule \textsc{Abs} and for Section~\ref{sec:safety}. The lambda
calculus is a formal system in which everything is a function. A LETREC
\texttt{proc (x) body} is a lambda abstraction $\lambda x.\, \mathit{body}$, and the
application \texttt{(f e)} is the application $f\ e$.

\subsection{Syntax}

There are three constructs: variables $x, y, z, \dots$; abstraction
$\lambda x.\, e$, a function with parameter $x$ and body $e$; and application
$e_1\, e_2$, which applies $e_1$ to $e_2$. Application associates to the left
($f\,a\,b$ is $(f\,a)\,b$), and the body of an abstraction extends as far to
the right as possible. Figure~\ref{fig:lambda-ast} shows the abstract syntax
tree of $(\lambda x.\, x\, x)\,(\lambda y.\, y)$.

\begin{figure}[h]
\centering
\begin{forest}
  for tree={font=\small, s sep=5mm, l sep=4mm},
  [app
    [$\lambda x$ [app [$x$] [$x$]]]
    [$\lambda y$ [$y$]]
  ]
\end{forest}
\caption{Abstract syntax tree of $(\lambda x.\, x\, x)\,(\lambda y.\, y)$.}
\label{fig:lambda-ast}
\end{figure}

\paragraph{Free and bound variables.} In $\lambda x.\, M$ the variable $x$ is
\textbf{bound}: the abstraction is its scope. A variable that is not inside an
abstraction for it is \textbf{free}. In $\lambda x.\, x\, y$, $x$ is bound and
$y$ is free. The name of a bound variable does not matter, as long as renaming
does not capture another variable: $\lambda x.\, x \equiv_\alpha \lambda y.\, y$
($\alpha$-equivalence).

\subsection{Reduction}

Computation is one rule, $\beta$-reduction, which substitutes the argument into
the body:
\[
(\lambda x.\, e)\, v \;\to_\beta\; e[x := v].
\]
A second rule, $\eta$-reduction, says that a function that only passes its
argument on is the function itself: $\lambda x.\, (f\, x) \to_\eta f$ when $x$ is
not free in $f$.

For $(\lambda x.\, x\, x)\,(\lambda y.\, y)$:
\[
(\lambda x.\, x\, x)\,(\lambda y.\, y) \;\to_\beta\; (\lambda y.\, y)\,(\lambda y.\, y)
\;\to_\beta\; \lambda y.\, y.
\]
The result has no more redexes (subterms of the form $(\lambda x.\, e)\, v$); it is
in normal form.

\paragraph{Evaluation strategies.} When a term has several redexes, the order
matters. Normal order reduces the leftmost, outermost redex first, and finds
a normal form whenever one exists. Call by value evaluates the argument before
the call, as C, Java and Python do. Call by name passes the argument
unevaluated, which is the basis of lazy evaluation. The runtime block
returns to these strategies with the lecture's examples.

\subsection{Church encodings}

The pure calculus has no numbers or booleans, but they can be encoded as
functions. A boolean chooses between two alternatives:
\[
\mathit{TRUE} \equiv \lambda t.\, \lambda f.\, t \qquad
\mathit{FALSE} \equiv \lambda t.\, \lambda f.\, f \qquad
\mathit{AND} \equiv \lambda p.\, \lambda q.\, p\, q\, p.
\]
A Church numeral $n$ applies a function $n$ times:
\[
0 \equiv \lambda f.\, \lambda x.\, x \qquad
1 \equiv \lambda f.\, \lambda x.\, f\, x \qquad
2 \equiv \lambda f.\, \lambda x.\, f\,(f\, x) \qquad
\mathit{SUCC} \equiv \lambda n.\, \lambda f.\, \lambda x.\, f\,(n\, f\, x).
\]

\subsection{Recursion}

Abstractions are anonymous, so a function cannot call itself by name. LETREC
solves this with \texttt{letrec}, which puts the function in its own scope.
The pure calculus uses a fixed-point combinator:
\[
Y = \lambda f.\, (\lambda x.\, f\,(x\, x))\,(\lambda x.\, f\,(x\, x)).
\]
One step of reduction shows that $Y\, g \to_\beta (\lambda x.\, g\,(x\, x))\,(\lambda x.\, g\,(x\, x)) \to_\beta g\,((\lambda x.\, g\,(x\, x))\,(\lambda x.\, g\,(x\, x)))$,
and the argument of $g$ is again the term that $Y\, g$ reduced to. So
$Y\, g$ behaves as $g\,(Y\, g)$: it gives $g$ a copy of itself.

\subsection{Exercises}

\begin{exercise}
Define $\mathit{OR}$ and $\mathit{NOT}$ with Church booleans and check that
$\mathit{OR}\ \mathit{TRUE}\ \mathit{FALSE}$ reduces to $\mathit{TRUE}$.
\end{exercise}

\paragraph{Solution.} $\mathit{OR} \equiv \lambda p.\, \lambda q.\, p\, p\, q$ and
$\mathit{NOT} \equiv \lambda p.\, p\, \mathit{FALSE}\, \mathit{TRUE}$. Then
$\mathit{OR}\ \mathit{TRUE}\ \mathit{FALSE} \to_\beta \mathit{TRUE}\ \mathit{TRUE}\ \mathit{FALSE}
\to_\beta \mathit{TRUE}$, because $\mathit{TRUE}$ returns its first argument.

\begin{exercise}
Show that $\mathit{SUCC}\ 1$ reduces to $2$.
\end{exercise}

\paragraph{Solution.} $\mathit{SUCC}\ 1 \to_\beta \lambda f.\, \lambda x.\, f\,(1\, f\, x)$,
and $1\, f\, x = (\lambda f.\, \lambda x.\, f\, x)\, f\, x \to_\beta f\, x$, so the result is
$\lambda f.\, \lambda x.\, f\,(f\, x) = 2$.

\begin{exercise}
Which variables are free in $\lambda x.\, (\lambda y.\, x\, y\, z)\, y$?
\end{exercise}

\paragraph{Solution.} $z$, and the last $y$, which is outside the abstraction
$\lambda y$. The first $y$ and $x$ are bound.
